% !PRINTED 548
\weberpage{548}
%<<CONT
in which $p_1, p_2 \ldots$ are powers of \gesperrt{different} primes.
We further set
\[
m = p_1 m_1 = p_2 m_2 \ldots,
\]
and choose, as is always possible by what has just been proved,
$\lambda_1$ relatively prime to $p_1$, $\lambda_2$ relatively prime to $p_2$ etc., so that
\[
(\varepsilon^{\lambda_1}, \alpha), \quad (\varepsilon^{\lambda_2}, \alpha) \ldots
\]
are different from zero. Then by the Theorem \S.\,164, 3.:
\weq{5}{(\varepsilon^{\lambda}, \alpha)(\varepsilon^{\lambda_1}, \alpha)^{m_1 \nu}(\varepsilon^{\lambda_2}, \alpha)^{m_2 \nu} \ldots = \chi_{\lambda}}
is a quantity contained in $\Omega(\varepsilon)$, if
\weq{6}{\lambda \equiv -\nu(\lambda_1 m_1 + \lambda_2 m_2 + \cdots) \ (\mathrm{mod}\ m).}
But $(\varepsilon^{\lambda_1}, \alpha)^{m_1}$ is a root of degree $p_1$ of a function
in $\Omega(\varepsilon)$, and we therefore set
% [sic] authority omits the exponent m_2 from the second factor; symmetry with the first factor requires it.
\weq{7}{(\varepsilon^{\lambda_1}, \alpha)^{m_1} = \sqrt[p_1]{\varphi_1}, \quad (\varepsilon^{\lambda_2}, \alpha) = \sqrt[p_2]{\varphi_2} \ldots, \quad (\varepsilon^{\lambda}, \alpha) = \sqrt[m]{\psi_{\lambda}}.}
Then by (5) we have
\weq{8}{\sqrt[m]{\psi_{\lambda}} = \frac{\chi_{\lambda}}{\left(\sqrt[p_1]{\varphi_1}\sqrt[p_2]{\varphi_2} \ldots \right)^{\nu}}.}
Now $(\lambda_1 m_1 + \lambda_2 m_2 + \cdots)$ is relatively prime to $m$, since $m_2 \ldots$
are divisible by $p_1$, while $\lambda_1 m_1$ is relatively prime to $p_1$; and thus from (6)
one obtains for each $\lambda$ a number $\nu$ completely determined
with respect to the modulus $m$.
If, then, we substitute the expressions (8) in (3) and assign to the
radicals $\sqrt[p_1]{\varphi_1}, \sqrt[p_2]{\varphi_2} \ldots$ all their values, we obtain
exactly $m$ different values for $\alpha$, and no more.

The last results we can use to derive a form of the
representation of the roots $\alpha$ in a somewhat generalised shape,
which \gesperrt{Abel} communicates in the place cited, and which
relates to the case where the field $\Omega$ is real, i.e. consists of
nothing but real numbers. The functions $\varphi, \psi, \chi$, as we
have used them above, are then composed of real
numbers and of the root of unity $\varepsilon$, which one finds by the division
of the circumference of the circle into $m$ equal parts; one may, for instance, set
\[
\varepsilon = e^{\frac{2\pi i}{m}} = \cos\frac{2\pi}{m} + i \sin\frac{2\pi}{m}
\]

The function $\varphi_1$ passes over, when $\varepsilon$ is changed into $\varepsilon^{-1}$, into
the conjugate imaginary value, which we denote by $\varphi_1'$.
